% =====================================================================
%  Comparison with the concurrent convergence claim of Mishra, Trivedi
%  and Kumar (2026) for SignMuon.  Drop-in appendix subsection for
%  the one-column appendix builds.  Self-contained: refers only to this
%  paper's Theorem~\ref{th:1}, Proposition~\ref{prop:reduction}, the
%  linear objective~\eqref{eq:lin_obj}, Remark~\ref{rem:efsm_bdd}, and
%  \citet{mishra2026signmuon}.
% =====================================================================
\subsection{Comparison with the convergence claim for SignMuon}\label{app:mishra}

Concurrently, \citet{mishra2026signmuon} introduced the sign-after-LMO method
under the same name \emph{SignMuon}: their algorithm forms the momentum
$\mathbf{M}_t$, computes its polar factor, and steps along
$\mathbf{S}_t=\operatorname{sign}(\operatorname{polar}(\mathbf{M}_t))$,
normalized to $\mathbf{D}_t=\mathbf{S}_t/\sqrt{mn}$; after absorbing
$1/\sqrt{mn}$ into $\eta$ this is the SignMuon step of
Section~\ref{sec:theory}. Their abstract and contributions attribute to this
method an $\mathcal{O}(1/\sqrt{T})$ stationarity guarantee; the theorem that
establishes the rate is stated, accurately, for a \emph{gradient-sign
instantiation}. The distance between the two statements is the subject of this
subsection. Two facts resolve it, and neither contradicts
Theorem~\ref{th:1}: the finalized rate is proved for an update that computes
no polar factor and carries no momentum, and the one bound of theirs that does
apply to the SignMuon update is an inequality whose right-hand side exceeds
its left-hand side on the instance of Theorem~\ref{th:1}, at every iteration
and for every step size, so that it is satisfied there while the method
ascends.

\paragraph{The two components of their analysis.}
Their stationarity measure is
$\mathcal{G}_T=\tfrac1T\sum_{t=0}^{T-1}\mathbb{E}\bigl[\|\nabla
f(\mathbf{X}_t)\|_1/\sqrt{mn}\bigr]$, controlled through the per-entry
sign-error probabilities
$q_{ij,t}=\Pr\bigl(\mathbf{S}_{t,ij}\neq\operatorname{sign}([\nabla
f(\mathbf{X}_t)]_{ij})\mid\mathbf{X}_t\bigr)$ of the transmitted sign matrix
$\mathbf{S}_t$. The first component is generic. For an \emph{arbitrary} sign
oracle, the conditional identity
\begin{equation}\label{eq:mishra_identity}
\mathbb{E}\bigl[\langle \nabla f(\mathbf{X}_t),\,\mathbf{D}_t\rangle \,\big|\, \mathbf{X}_t\bigr]
=\frac{1}{\sqrt{mn}}\sum_{i,j}\bigl|[\nabla f(\mathbf{X}_t)]_{ij}\bigr|\,\bigl(1-2q_{ij,t}\bigr)
\end{equation}
and the descent lemma of spectral smoothness telescope into
\begin{equation}\label{eq:mishra_generic}
\mathcal{G}_T
\;\le\;
\frac{f(\mathbf{X}_0)-f^\ast}{\eta T}
+\frac{L_\ast}{2}\,\eta
+\underbrace{\frac{2}{T}\sum_{t=0}^{T-1}\mathbb{E}\!\left[\frac{1}{\sqrt{mn}}\sum_{i,j}\bigl|[\nabla f(\mathbf{X}_t)]_{ij}\bigr|\,q_{ij,t}\right]}_{=:\,R_T}.
\end{equation}
The second component estimates the residual $R_T$, and it is here that the
oracle is fixed: for $\mathbf{S}_t=\operatorname{sign}(\widetilde{\mathbf{G}}_t)$,
where $\widetilde{\mathbf{G}}_t$ is an unbiased stochastic gradient with
coordinatewise variance at most $\sigma_{ij}^2/n_b$ at mini-batch size $n_b$,
a Markov--Jensen argument yields $\bigl|[\nabla
f(\mathbf{X}_t)]_{ij}\bigr|\,q_{ij,t}\le\sigma_{ij}/\sqrt{n_b}$, hence
$R_T\le2\|\sigma\|_1/\sqrt{mn\,n_b}$ with
$\|\sigma\|_1=\sum_{i,j}\sigma_{ij}$. The residual vanishes as
$n_b\to\infty$, and the choice $n_b=T$ produces the
$\mathcal{O}(1/\sqrt{T})$ rate.

\paragraph{The finalized rate is not a result about SignMuon.}
The oracle of the second component transmits the sign of the stochastic
gradient itself. That update invokes neither the momentum buffer nor the polar
factor; as an algorithm it is SignSGD at batch size $n_b$, in single-worker
and majority-vote form, normalized by $1/\sqrt{mn}$ and analysed under
spectral rather than coordinatewise smoothness, which their own comparison
identifies as the sole improvement over \citet{bernstein2018signsgd}. Nor is
the restriction incidental. The Markov--Jensen estimate bounds the probability
of a sign error by the ratio of noise to signal,
$\sigma_{ij}/(\sqrt{n_b}\,|[\nabla f(\mathbf{X}_t)]_{ij}|)$, and is therefore
available exactly when sampling noise is the only mechanism by which a
transmitted sign can disagree with the gradient's. For the SignMuon oracle the
disagreement is structural rather than stochastic: even with exact gradients
($\sigma\equiv0$), $q_{ij,t}$ is the indicator that
$\operatorname{sign}(\operatorname{polar}(\mathbf{M}_t))$ and
$\operatorname{sign}(\nabla f(\mathbf{X}_t))$ differ at $(i,j)$, a quantity
that no batch size reduces. The $\mathcal{O}(1/\sqrt{T})$ rate accordingly
attaches to the gradient-sign update, and to SignMuon their analysis offers
only \eqref{eq:mishra_generic} with $R_T$ unestimated.

\paragraph{Scope of the generic bound.}
Inequality \eqref{eq:mishra_generic} is valid for every sign oracle, and for
that reason asserts nothing until $R_T$ is estimated. By
\eqref{eq:mishra_identity}, the summand of $R_T$ at time $t$ exceeds the
corresponding summand of $\mathcal{G}_T$ precisely when
$\mathbb{E}[\langle\nabla f(\mathbf{X}_t),\mathbf{D}_t\rangle\mid\mathbf{X}_t]\le0$,
that is, precisely when the expected step fails to be a descent direction.
Whenever this occurs at every $t$, the right-hand side of
\eqref{eq:mishra_generic} exceeds the left-hand side termwise and the
inequality holds irrespective of how the iterates behave. The bound therefore
has content only where the transmitted sign is already positively aligned with
the gradient in expectation; that alignment is the property a convergence
proof for SignMuon would have to establish, and it is the property
Theorem~\ref{th:1} refutes.

\paragraph{On the instance of Theorem~\ref{th:1}.}
Consider the linear objective~\eqref{eq:lin_obj} with the $4\times4$ gradient
$\mathbf{G}$ of Theorem~\ref{th:1}. By Proposition~\ref{prop:reduction} the
gradient equals $\mathbf{G}$ and the update direction equals the constant
matrix $\operatorname{sign}(\operatorname{polar}(\mathbf{G}))$ at every
iteration, whatever the momentum, so the oracle is deterministic and $q_{ij}$
is the indicator that
$\operatorname{sign}(\operatorname{polar}(\mathbf{G}))_{ij}\neq
\operatorname{sign}(\mathbf{G}_{ij})$. Identity \eqref{eq:mishra_identity}
then evaluates exactly:
\begin{equation}\label{eq:mishra_instance}
\sum_{i,j}|\mathbf{G}_{ij}|\,(1-2q_{ij})
=\bigl\langle \mathbf{G},\,\operatorname{sign}(\operatorname{polar}(\mathbf{G}))\bigr\rangle
=-\tfrac{42468}{103}\;<\;0 .
\end{equation}
Dividing by $\sqrt{mn}=4$, every term of $R_T$ exceeds the corresponding term
of $\mathcal{G}_T$ by the same amount, so that
\begin{equation}\label{eq:mishra_dominates}
R_T=\mathcal{G}_T+\tfrac{10617}{103}
\qquad\text{for every } T .
\end{equation}
The right-hand side of \eqref{eq:mishra_generic} therefore exceeds the
left-hand side by at least $10617/103$ for every $\eta>0$, every $L_\ast$ and
every $T$: the inequality is satisfied and constrains nothing. What the
trajectory actually does is read off the same identity,
$f(\mathbf{X}_t)-f(\mathbf{X}_{t-1})=-\eta\langle\mathbf{G},\mathbf{D}_t\rangle
=\eta\cdot\tfrac{10617}{103}>0$, the divergence of Theorem~\ref{th:1}: the
excess $R_T-\mathcal{G}_T$ of \eqref{eq:mishra_dominates} and the
per-iteration ascent rate
$\eta^{-1}\bigl(f(\mathbf{X}_t)-f(\mathbf{X}_{t-1})\bigr)$ are the same
number.
The one assumption of theirs the linear instance lacks is lower boundedness,
and the modification of Remark~\ref{rem:efsm_bdd} applies unchanged: $f$ is
unbounded below only on a half-space the iterates never enter, where a smooth
bounded replacement restores the assumption without moving the trajectory or
any quantity above.

In their terms, Theorem~\ref{th:1} exhibits a smooth instance on
which the sign-error probabilities $q_{ij,t}$ of the SignMuon oracle, averaged
over the entries with weights $\bigl|[\nabla f(\mathbf{X}_t)]_{ij}\bigr|$,
exceed $\tfrac12$ at every iteration. Any rate extracted from
\eqref{eq:mishra_generic} requires that weighted average to stay below
$\tfrac12$ by a uniform margin, and no assumption of theirs implies such a bound for
$\operatorname{sign}(\operatorname{polar}(\mathbf{M}_t))$. Their theorems
stand as guarantees for majority-vote SignSGD under spectral smoothness. A
convergence guarantee for SignMuon they are not, and Theorem~\ref{th:1} shows
that none is available at this level of generality.
